% Source PDF pages 41--42; manual pages 2-12--2-13.
\subsection{Local Geodetic Horizon Coordinate System}\label{sec:local-geodetic}

The local geodetic horizon coordinate system is similar to the local geocentric
horizon system, except that its horizon plane is referenced to an ellipsoid rather
than to a sphere. The major axes of this ellipsoid of revolution lie in the equatorial
plane, as shown in \cref{fig:local-geodetic-horizon}. Its shape is controlled by the
equatorial earth radius $R_{\oplus}$ and any one of three redundant quantities: the
polar radius $R_p$, eccentricity $e$, or flattening $f$. They are related by
\begin{equation}
  R_p = R_{\oplus}\sqrt{1-e^2}=R_{\oplus}(1-f).
  \label{eq:ellipsoid-parameter-relation}
\end{equation}
The eccentricity and flattening must be less than one, and the polar radius must be
less than the equatorial radius. These parameters normally conform to WGS-72\taosref{5}
or WGS-84\taosref{6}, but they can be changed through the earth-model input described
in \cref{sec:block-earth}.

\begin{figure}[htbp]
  \centering
  \resizebox{0.72\linewidth}{!}{\begin{tikzpicture}[scale=0.95,>=Latex,line cap=round,line join=round]
  \coordinate (O) at (0,0);
  \coordinate (V) at (4.25,2.45);
  \coordinate (S) at (3.77,2.05);
  % Ellipsoidal earth section.
  \draw[line width=0.95pt] (0,0) arc[start angle=180,end angle=0,x radius=3.75,y radius=2.30];
  \draw[->,line width=0.9pt] (O) -- (7.30,0) node[right,align=left] {Equatorial\\plane};
  \draw[->,line width=0.9pt] (O) -- (0,3.65) node[above] {$\hat{z}_{\oplus}$};
  \draw[<->] (-0.40,0) -- (-0.40,2.30) node[midway,left] {$R_p$};
  \draw[<->] (0,-0.36) -- (3.75,-0.36) node[midway,below] {$R_{\oplus}$};
  \draw[->,line width=1pt] (O) -- (V) node[midway,above left] {$\vec r$};
  \fill (V) circle (1.5pt);
  \fill (S) circle (1.3pt);
  % Surface normal and tangent plane.
  \draw[line width=0.9pt] (2.35,0.33) -- (V);
  \draw[line width=0.9pt] ($(V)+(-1.25,0.95)$) -- ($(V)+(1.20,-0.90)$);
  \node[align=left,above] at (4.25,3.08) {Geodetic local\\horizon plane};
  \draw[->,line width=1pt] (V) -- ($(V)+(-0.92,0.68)$) node[above left] {$\hat{x}_{gd}$};
  \draw[->,line width=1pt] (V) -- ($(V)+(0.25,0.78)$) node[above] {$\hat{y}_{gd}$};
  \draw[->,line width=1pt] (V) -- ($(V)+(-0.66,-0.53)$) node[below left] {$\hat{z}_{gd}$};
  \draw[->] (0.92,0) arc[start angle=0,end angle=31,radius=0.92];
  \node at (0.92,0.34) {$\delta_{gc}$};
  \draw[->] (2.45,0) arc[start angle=0,end angle=47,radius=0.72];
  \node at (2.50,0.42) {$\delta_{gd}$};
  \draw (3.58,1.80) -- ++(0.18,0.18) -- ++(-0.18,0.18);
\end{tikzpicture}
}
  \caption{Local geodetic horizon coordinates.}
  \label{fig:local-geodetic-horizon}
\end{figure}

The local geodetic horizon system has its origin at the vehicle's center of mass. Its
$z$ axis points downward, perpendicular to the ellipsoidal surface. The local horizon
plane is normal to this $z$ axis and therefore parallel to the plane tangent to the
earth's surface below the vehicle.

The $x$ and $y$ axes lie in the local horizon plane. The $x$ axis is in the plane
formed by $\vec{r}$ and $\hat{z}_{\oplus}$ and points north. The $y$ axis is
perpendicular to the $x$ axis and points east.

The local geodetic and local geocentric systems have the same origin at the vehicle's
center of mass. Their east vectors are identical, $\hat{y}_{gc}=\hat{y}_{gd}$, while
the north and down vectors are rotated about the common $y$ axis by
$\delta_{gd}-\delta_{gc}$.

The geodetic unit vectors are
\begin{equation}
  \hat{x}_{gd}
    = -\sin\delta_{gd}
       (\cos\lambda\,\hat{x}_{\oplus}
        +\sin\lambda\,\hat{y}_{\oplus})
       +\cos\delta_{gd}\,\hat{z}_{\oplus},
  \label{eq:local-geodetic-x-unit-vector}
\end{equation}
\begin{equation}
  \hat{y}_{gd}
    = -\sin\lambda\,\hat{x}_{\oplus}
      +\cos\lambda\,\hat{y}_{\oplus},
  \label{eq:local-geodetic-y-unit-vector}
\end{equation}
\begin{equation}
  \hat{z}_{gd}
    = -\cos\delta_{gd}
       (\cos\lambda\,\hat{x}_{\oplus}
        +\sin\lambda\,\hat{y}_{\oplus})
      -\sin\delta_{gd}\,\hat{z}_{\oplus}.
  \label{eq:local-geodetic-z-unit-vector}
\end{equation}
These equations are identical in form to
\cref{eq:local-geocentric-unit-vectors}, with geodetic latitude replacing
geocentric latitude. The latitude required by these equations is obtained with the
iterative method in \cref{sec:geodetic}, and the vectors are computed by
\texttt{geod\_unit\_vectors}.
